\chapter{Conclusion generale}

Ce stage a eu pour objectif de concevoir, realiser et consolider une plateforme de gestion intelligente des tickets support capable de repondre a des besoins metier reels : centralisation des demandes, clarification des responsabilites, supervision des delais, archivage documentaire, assistance a la decision et preparation d'une demonstration techniquement defendable.

Le projet SupportFlow montre qu'une application de support ne peut pas etre reduite a une simple logique de creation et de cloture de tickets. Pour etre utile, elle doit prendre en compte les acteurs, les regles metier, les contraintes de securite, le suivi SLA, la capitalisation documentaire et les exigences d'observabilite. C'est pourquoi la solution proposee s'appuie sur une architecture riche, combinant Angular, Spring Boot, Camunda, Keycloak, Alfresco, Docker, SonarQube, GitHub Actions, ArgoCD et une brique IA locale.

Au fil du stage, plusieurs dimensions ont ete traitees : l'analyse du besoin, la conception de l'architecture, la realisation technique, la gestion des roles et des permissions, la structuration des parcours client, agent et manager, la reprise des workflows Camunda, la stabilisation de la chaine GitOps, l'integration de la GED, la generation des rapports et la mise en place de preuves qualite. Cette progression a permis de faire evoluer SupportFlow d'un ensemble de fonctionnalites vers un systeme metier plus coherent, plus explicable et plus solide.

Les resultats obtenus sont significatifs. L'application permet aujourd'hui de gerer des tickets selon un cycle de vie structure, d'exposer des parcours differencies selon les roles, de surveiller les tickets a risque, de produire des rapports, d'archiver les documents, d'utiliser un copilote IA et de s'inscrire dans une demarche DevOps et GitOps documentee. Le projet ne se contente donc pas de fonctionner localement : il dispose aussi d'un socle de qualite, de deploiement et de soutenance.

Sur le plan personnel, ce stage a ete formateur a plusieurs niveaux. Il a permis de renforcer les competences techniques en developpement full stack, securite, BPMN, conteneurisation et integration continue. Il a aussi developpe des competences methodologiques essentielles : analyse du besoin, priorisation, gestion des risques, correction transverse, documentation et preparation d'une demonstration convaincante.

Bien entendu, le projet n'est pas fige. Plusieurs perspectives existent : renforcer encore les tests de non-regression, enrichir la base de connaissances, pousser plus loin les vues manager de supervision, rendre l'IA plus observable et industrialiser davantage le deploiement multi-environnements. Ces limites ne diminuent pas la valeur du travail realise ; elles montrent au contraire que SupportFlow constitue une base serieuse pour des evolutions futures.

En conclusion, le stage a permis de repondre de maniere satisfaisante a la problematique initiale. SupportFlow apporte une solution credible, integree et evolutive pour la gestion des tickets support. Il illustre la capacite a conduire un projet complexe de bout en bout, depuis le besoin metier jusqu'a la preuve technique, tout en produisant un resultat coherent pour l'entreprise et pertinent dans un cadre academique.
